\chapter{پاسخ مسئله‌ها}
\thispagestyle{empty}

\begin{quote}
	{\lotus
		اگر مسئله‌ای وجود دارد که نمی‌توانید حل کنید، مسئله‌ی ساده‌تری وجود دارد که می‌توانید حل کنید؛ آن را بیابید. \quad - پولیا%
	}
\end{quote}
%%%%%%%%%%%%%%%%%%%%%%%%%%%%%%%%%%%%%%%%%%%%%%%
\section*{منطق ریاضی}
\textbf{\ref{loprpedar}}.
در نوبت اول هر دو هم‌زمان می‌گویند: «خیر.» چون هر کدام این‌طور فکر می‌کنند که پیشانی دیگری گلی است و حداقل پیشانی یکی از ما گلی است، پس نمی‌توانم بفهمم پیشانی من گلی است یا نه. در نوبت دوم هر دو هم‌زمان می‌گویند: «بله.» چون بعد از پاسخ قبلی، هر کدام این‌طور فکر می‌کنند که اگر پیشانی من گلی نبود، دیگری باید پاسخ مثبت می‌داد، پس پیشانی من گلی است.

\df \\ \textbf{\ref{loprabccolor}}.
فرض کنید تمبر C قرمز باشد. A پاسخ منفی داده، پس تمبر B قرمز نبوده است. اما اگر چنین است، B نباید پاسخ منفی می‌داد چون از روی پاسخ منفی A و رنگ تمبر C متوجه می‌شد که رنگ تمبرش قرمز نیست. پس تمبر C نمی‌تواند قرمز باشد. با استدلال مشابه (چون تعداد تمبرهای قرمز و زرد یکسان است)، می‌توان گفت که تمبر C نمی‌تواند زرد باشد. بنابراین تمبر C سبز است. اما درباره‌ی رنگ تمبرهای A و B نمی‌توان نظر داد، ممکن است هر کدام از سه رنگ باشند.

\df \\ \textbf{\ref{loprmcgregor}}.
خانه‌ی اول: اگر مرد شوالیه بود، هیچ وقت ادعا نمی‌کرد که او و همسرش هر دو سربازند. در نتیجه، او باید سرباز باشد. از آن‌جا که او سرباز است، جمله‌اش نادرست است؛ پس آن‌ها هر دو سرباز نیستند. این یعنی همسرش باید شوالیه باشد. \\
خانه‌ی دوم: اگر مرد سرباز بود، پس این‌که حداقل یکی از آن‌ها سرباز است، درست می‌شد، اما ممکن نیست، چون سرباز نمی‌تواند راست بگوید. پس مرد شوالیه است. این یعنی ادعایش درست است و حداقل یکی از آن دو سرباز است. از آن‌جا که خودش سرباز نیست، باید همسرش سرباز باشد. \\
خانه‌ی سوم: فرض کنید مرد شوالیه باشد. پس آنچه گفته درست است، پس همسرش باید شوالیه باشد. اگر مرد شوالیه باشد، همسرش هم شوالیه است. این جمله درست است. پس مرد شوالیه است و بنابراین همسرش نیز شوالیه است. \\
خانه‌ی چهارم: نمی‌توان مشخص کرد که مرد شوالیه است یا سرباز، اما نوع همسرش را می‌توان به این صورت مشخص کرد: اگر زن سرباز بود، شوهرش هیچ وقت ادعا نمی‌کرد که نوع او با همسرش یکسان است، چون اگر مرد سرباز باشد، جمله‌ی درست نمی‌گوید. پس زن باید شوالیه باشد.

\df \\ \textbf{\ref{loprladytiger}}.
دو حالت وجود دارد: \\
حالت یکم: علامت در ۱ درست و علامت در ۲ نادرست است. درستی علامت در ۱ نشان می‌دهد که در اتاق ۱ یک دختر و در اتاق ۲ یک ببر است. نادرستی علامت در ۲ نشان می‌دهد که در هر دو اتاق دختر یا در هر دو اتاق ببر است. اما دو نتیجه با هم تناقض دارند، پس حالت یکم ممکن نیست. \\
حالت دوم:  علامت در ۱ نادرست و علامت در ۲ درست است. نادرستی علامت در ۱ نشان می‌دهد که در اتاق ۱ یک ببر و در اتاق ۲ یک دختر است. درستی علامت ۲ نشان می‌دهد که در یکی از دو اتاق دختر و دیگری ببر است که با نتیجه‌ی نادرستی علامت در ۱ منطبق است. پس حالت دوم ممکن است. بنابراین، زندانی باید در ۲ را باز کند.

\df \\ \textbf{\ref{loprpirates}}.
فرض کنید نوشته‌ی صندوقچه‌ی ۲ درست باشد، یعنی صندوقچه‌ی ۲ خالی است. نوشته‌ی صندوقچه‌ی ۱ نادرست است، یعنی صندوقچه‌ی ۱ خالی نیست. نوشته‌ی صندوقچه‌ی ۳ نادرست است، یعنی گنج در صندوقچه‌ی ۲ نیست. بنابراین، گنج در صندوقچه‌ی ۱ است.

\df \\ \textbf{\ref{loprkraig}}.
فرض می‌کنیم دو عبارت X و Y داشته باشیم که یا هر دو درست یا هر دو نادرست باشند. در این صورت اگر یکی از افراد تیمارستان یکی از این عبارت‌ها را باور داشته باشد، عبارت دیگر را نیز باور خواهد داشت. اگر هر دو عبارت درست باشد، در آن صورت هر کدام از افراد تیمارستان که یکی 	از آن‌ها را باور داشته باشد، بایستی عاقل باشد و بنابراین باید عبارت دیگر را باور داشته باشد. اگر هر دو عبارت نادرست باشد، هرکس که یکی از آن‌ها را باور داشته باشد، بایستی دیوانه باشد و باید عبارت دیگر را باور داشته باشد.

دو کمیته‌ی دلخواه مثلاً کمیته‌ی B و C را در نظر می‌گیریم. مجموعه‌ی افرادی را که بدترین دشمنشان در کمیته‌ی B است، U و مجموعه‌ی افرادی را که بهترین دوستشان در کمیته‌ی C است، V می‌نامیم. بنا به مورد ۴، مجموعه‌های U و V نیز هر دو کمیته هستند. بنا به مورد ۵، فردی وجود دارد مثلاً به نام دان که بهترین دوست او مثلاً به نام ادوارد عقیده دارد که او در کمیته‌ی U است و بدترین دشمن او مثلاً به نام فرد عقیده دارد که او در کمیته‌ی V است. اکنون طبق تعریف U اگر بگوییم دان در کمیته‌ی U است، مثل آن است که بگوییم بدترین دشمن او یعنی فرد در کمیته‌ی B است. به بیان دیگر، دو عبارت «دان در U است» و «فرد در کمیته‌ی B است» یا هر دو درست یا هر دو نادرست هستند. چون ادوارد یکی از این دو عبارت یعنی بودن دان در U را باور دارد، پس بایستی عبارت دیگر یعنی بودن فرد در کمیته‌ی B را نیز باور داشته باشد. پس ادوارد عقیده دارد که فرد در کمیته‌ی B است.

از سوی دیگر، فرد عقیده دارد که دان در کمیته‌ی V است. طبق تعریف \lr{V}، دان فقط در صورتی عضو V خواهد بود که دوست او ادوارد عضو کمیته‌ی C باشد، یعنی این دو عبارت یا هر دو درست یا هر دو نادرست هستند. اما چون فرد عقیده دارد که دان عضو V است، فرد بایستی عقیده داشته باشد که ادوارد عضو کمیته‌ی C است.

اکنون کمیته‌ی بیماران و کمیته‌ی پزشکان را مطابق موردهای ۱ و ۲ در نظر بگیرید و این دو کمیته را به ترتیب به‌جای کمیته‌های B و C فرض کنید. ‌دو نفر مانند ادوارد و فرد وجود دارند که دارای عقیده‌های زیر هستند: ادوارد عقیده دارد که فرد در کمیته‌ی بیماران است و فرد عقیده دارد که ادوارد در کمیته‌ی پزشکان است.

فرد یا عاقل است یا دیوانه. فرض کنیم عاقل باشد. در این صورت، باور او درست است، یعنی ادوارد واقعاً پزشک است. اگر ادوارد دیوانه باشد، در این صورت پزشکی دیوانه است. اگر ادوارد عاقل باشد، در این صورت باور او درست است، یعنی فرد بیمار و بنا به فرض، بیماری عاقل است. به این ترتیب ثابت می‌شود که اگر فرد عاقل باشد، یا او بیماری عاقل یا ادوارد پزشکی دیوانه است.

اکنون فرض کنیم که فرد دیوانه باشد. در این صورت، باور او نادرست است، یعنی ادوارد بیمار است. اگر ادوارد بیمار باشد، در این صورت بیماری عاقل است. اگر ادوارد دیوانه باشد، در این صورت باور او نادرست است و نتیجه می‌شود که فرد پزشک و بنا به فرض، پزشکی دیوانه است. به این ترتیب ثابت می‌شود که اگر فرد دیوانه باشد، یا او پزشکی دیوانه است یا ادوارد بیماری عاقل است.

بنابراین، یکی از این دو نفر، یعنی ادوارد یا فرد، بایستی یا پزشکی دیوانه یا بیماری عاقل باشد.

\df \\ \textbf{\ref{lopr4cab45}}.
حداقل باید دو کارت $B$ و ۴ را برگردانیم. کارت ۵ تأثیری ندارد چون عددی فرد است. کارت $A$ نیز تأثیری ندارد چون حرف صدادار است.

\df \\ \textbf{\ref{loprrttp}}.
\begin{table}[H]
	\begin{latin}
		\centering
		\begin{tabular}{c|c|c|c}
			$p$ & $q$ & $p \vee q$ & $(p \vee q) \rightarrow q$ \\
			\hline
			T & T & T & T \\
			T & F & T & F \\
			F & T & T & T \\
			F & F & F & T 
		\end{tabular}
		\qquad \qquad
		\begin{tabular}{c|c|c|c|c}
			$p$ & $q$ & $\neg p$ & $p \leftrightarrow q$ & $\neg p \wedge (p \leftrightarrow q)$ \\
			\hline
			T & T & F & T & F \\
			T & F & F & F & F \\
			F & T & T & T & T \\
			F & F & T & T & T
		\end{tabular}
	\end{latin}
\end{table}

\df \\ \textbf{\ref{loprpqreqx}}.
دو روش برای نوشتن گزاره‌ی معادل وجود دارد. در روش اول سطرهایی را که ارزش گزاره‌ی $X$ درست است در نظر می‌گیریم. در روش دوم سطرهایی را که ارزش گزاره‌ی $X$ نادرست است در نظر می‌گیریم. \\
گزاره‌ی اول:
$$(p \wedge \neg q \wedge r) \vee (\neg p \wedge q \wedge \neg r) \vee (\neg p \wedge \neg q \wedge q)$$
گزاره‌ی دوم:
$$(\neg p \vee \neg q \vee \neg r) \wedge (\neg p \vee \neg q \vee r) \wedge (\neg p \vee q \vee r) \wedge (p \vee q \vee r)$$
%%%%%%%%%%%%%%%%%%%%%%%%%%%%%%%%%%%%%%%%%%%%%%
\section*{اعداد}
\textbf{\ref{nupr1000prnu}}.
عددهای
$$1001!+2 \ , \ 1001!+3 \ , \ \cdots \ , \ 1001!+1000 \ , \ 1001!+1001$$
هزار عدد متوالی‌اند که هیچ‌کدام اول نیستند.
%%%%%%%%%%%%%%%%%%%%%%%%%%%%%%%%%%%%%%%%%%%%%%%
\section*{استدلال و برهان}
%%%%%%%%%%%%%%%%%%%%%%%%%%%%%%%%%%%%%%%%%%%%%%%
\section*{نظریه‌ی مجموعه‌ها}
%%%%%%%%%%%%%%%%%%%%%%%%%%%%%%%%%%%%%%%%%%%%%%%
\section*{معادلات و نامعادلات}
\textbf{\ref{arb+bra}}.
با توجه به فرمول مجموع و حاصل‌ضرب جواب‌های معادله‌ی درجه‌ی دوم داریم:
$$\alpha + \beta = \frac{-(-12)}{1} = 12 \ , \quad \alpha \beta = \frac{4}{1} = 4$$
بنابراین:
$$(\alpha \sqrt{\beta} + \beta \sqrt{\alpha})^2 = \alpha^2\beta + \beta^2\alpha + 2\alpha\beta\sqrt{\alpha\beta} = \alpha\beta(\alpha+\beta)+2\alpha\beta\sqrt{\alpha\beta} = 4(12)+2(4)(\sqrt{4}) = 64$$
خواسته‌ی مسئله برابر ریشه‌ی دوم عبارت بالاست، اما باید در نظر بگیریم که ریشه‌ی دوم ۶۴ دو جواب دارد: ۸ و $-8$. با توجه به این‌که مجموع و حاصل‌ضرب ریشه‌های عبارت درجه‌ی دوم، هر دو مثبت هستند، پس هر دو ریشه مثبت است و در نتیجه مقدار $-8$ قابل قبول نیست. پس جواب مسئله برابر با ۸ است.
%%%%%%%%%%%%%%%%%%%%%%%%%%%%%%%%%%%%%%%%%%%%%%%
\section*{ماتریس و دترمینان}

\df \\ \textbf{\ref{maprabprob}}.
داریم:
$A+B=
\begin{bmatrix}
	a & b-7 \\
	-1 & 1
\end{bmatrix}$.
شرط وارون‌پذیری این است که
$|A+B| \neq 0$،
بنابراین:
$$|A+B|=a+b-7 \neq 0 \Rightarrow a+b \neq 7$$
بنابراین حالاتی از
$(a,b)$
که مجموع دو تاس برابر ۷ است، قابل قبول نیست، یعنی:
$$\text{حالات نامطلوب}={(1,6),(2,5),(3,4),(4,3),(5,2),(6,1)}$$
$$\Rightarrow n(\text{مطلوب}) =36-6=30 \Rightarrow P(\text{مطلوب})=\frac{30}{36}=\frac{5}{6}$$

\df \\ \textbf{\ref{maprr3adet}}.
از دو طرف تساوی، دترمینان می‌گیریم.
\[
\left| \sqrt{3}A \right| =
\begin{vmatrix}
	-2 & |A|+1 \\
	3|A| & -|A|
\end{vmatrix}
\Rightarrow
\left( \sqrt{3} \right)^2 |A| = 2|A| - 3|A|^2 -3|A|
\Rightarrow
3|A|^2+4|A|=0
\Rightarrow
\begin{cases}
	|A|=0 \\
	|A|=-\frac{4}{3}
\end{cases}
\]
از آن‌جا که $A$ وارون‌پذیر است،
$|A| \neq 0$.
پس تنها جواب قابل قبول
$|A| = -\frac{4}{3}$
است.

\df
%%%%%%%%%%%%%%%%%%%%%%%%%%%%%%%%%%%%%%%%%%%%%%%
